\subsection*{Independent structural audit of the mixed Ferrari cover}

\paragraph{FC.1 (Proved): the involution kernel and its exact derived subgroup.}
Use the literal generators of FP.26:
\[
 L=\widehat B_4/\overline{\langle z_4\rangle},\quad
 z_4=(abc)^4,\quad z_3=(ab)^3,\quad
 x=ca^{-1},\quad y=aba c^{-1}b^{-1}a^{-1}.
\]
The fold sends $a,c$ to $a$ and $b$ to $b$. Its free kernel
before quotienting the centres is $K=\widehat F(x,y)$, with
the exact Nielsen actions
\[
 a:(x,y)\mapsto(x,yx),\qquad
 b:(x,y)\mapsto(xy^{-1},y).
 \tag{FC.1a}
\]
These are established directly from the Artin presentation in FP.26;
no geometric identification is needed for the group calculation below.
The fold of $z_4$ is $z_3^2$. Since $z_3$ has infinite
profinite order, the closed central subgroup generated by $z_4$
meets $K$ trivially. Thus $K$ embeds in $L$, with quotient
$\widehat B_3/\overline{\langle z_3^2\rangle}$.

Put $d=abc$ and $j=d^2$. The identity $c=xa$ and
$(ab)x(ab)^{-1}=y^{-1}$ give $d=y^{-1}(aba)$.
Using (FC.1a) once more gives
\[
       j=x^{-1}y^{-1}z_3,\qquad
       j^2=1,\qquad jxj=x^{-1},\quad jyj=y^{-1}.
 \tag{FC.1b}
\]
Consequently the kernel $C$ of the further quotient
$L\to\widehat B_3/\overline{\langle z_3\rangle}$ is
$K\rtimes\langle j\rangle$. The three elements
\[
             t_0=j,\qquad t_1=xj,\qquad t_2=yj
\]
are involutions. The inverse substitutions
$x=t_1t_0$, $y=t_2t_0$, $j=t_0$ identify the presentations,
so
\[
              C=C_2\amalg C_2\amalg C_2.
 \tag{FC.1c}
\]
The substitutions are valid as continuous maps of the indicated
profinite presentations, as well as for the discrete groups.

There is the useful exact identity
\[
                    \boxed{\ C'=K'K^2.\ }
 \tag{FC.1d}
\]
Here and below all derived and power subgroups are closed.
Indeed $C/(K'K^2)$ is abelian: the action of $j$ on
$K/(K'K^2)\simeq(C_2)^2$ is trivial. Conversely $K'\subset C'$
and $[x,j]=x^2$, $[y,j]=y^2$, which prove the reverse inclusion.
The subgroup $C'$ is therefore an index-four subgroup of the
free group $K$ of rank two. It is free of rank five, and
\[
  C'/(C''(C')^3)\simeq\mathbb F_3^5,\qquad
  |C/(C''(C')^3)|=8\cdot3^5=1944.
 \tag{FC.1e}
\]

\paragraph{FC.2 (Proved): exact actions, cusp, and the rank-five decomposition.}
The actions of $a,b$ on $C$ are
\[
\begin{split}
 a:(t_0,t_1,t_2)&\mapsto(t_0t_1t_0,t_0,t_2),\\
 b:(t_0,t_1,t_2)&\mapsto(t_2,t_1,t_2t_0t_2).
\end{split}
 \tag{FC.2a}
\]
For example (FC.1b), the centrality of $z_3$ in the
three-strand section, and (FC.1a) give
$aja^{-1}=x^{-1}j$ and $bjb^{-1}=yj$; the other four
images follow immediately. The generator $c=xa$ acts as
$\operatorname{Inn}(t_1t_0)$ followed by $a$.
Both $a$ and $b$ fix the peripheral element
\[
 w=t_1t_0t_2=xyj=z_3^{-1},\qquad w^2=[x,y].
 \tag{FC.2b}
\]
The last equality is a direct free-product reduction. In particular
the word order in $w$ is fixed here; permuting its three letters
would alter the marked cusp.

Write $Q=C/C'=(C_2)^3$, with the ordered basis $(t_0,t_1,t_2)$,
and $M=C'^{\rm ab}/3$. The Cayley covering graph for $C'$
is the cube, with eight vertices and twelve unoriented edges.
Its cellular modules over $\mathbb F_3$ are
\[
 C_0=\mathbb F_3[Q],\qquad
 C_1=\bigoplus_{i=0}^2
    \operatorname{Ind}_{\langle t_i\rangle}^{Q}\operatorname{sgn}.
\]
For a nontrivial character $\chi:Q\to\{\pm1\}$ of weight
$k$, the edge multiplicity is $k$, the vertex multiplicity is one,
and $H_0$ has multiplicity zero. Thus its multiplicity in $M$
is $k-1$. The trivial character occurs in neither $H_1$ nor any
edge summand. Hence
\[
 M=B\oplus P,\quad
 B=\chi_{110}\oplus\chi_{101}\oplus\chi_{011},\quad
 P=\chi_{111}^{\oplus2}.
 \tag{FC.2c}
\]
The even-parity subgroup $Q_0=K/C'$ is the kernel of
$\chi_{111}$. Since its order four is invertible in
$\mathbb F_3$, averaging over $Q_0$ is an idempotent projector
with image $P$. Transfer identifies this image with
$H_1(K,\mathbb F_3)$. Because $K\triangleleft L$ and
$C'\triangleleft L$, the splitting $M=B\oplus P$ is
$L$-invariant and canonical.

The independent program \texttt{check\_ferrari\_cube.py}
constructs the cube and a spanning tree, obtains five free
Schreier loops, and calculates every deck and $a,b$ matrix by
substituting the literal words (FC.2a). It imports no saved finite
quotient or matrix. It verifies the eight character multiplicities,
the two inverse automorphisms, the braid relation, (FC.1a), and
(FC.2b). The boundary linear image of $\langle a,b\rangle$
has order six, the proper linear image order twenty-four, and
their joint image order one hundred and forty-four.

The cusp class $[w^2]$ is a nonzero element of $B$ and is not
$Q$-invariant. This is verified by its explicit cube cycle, and
also follows from the geometric description below. Therefore the
central element $z_3^2=w^{-2}$ of the three-strand section acts
on $C/(C''(C')^3)$ by a nontrivial inner automorphism even though
its action on $M$ is trivial. Retaining only the rank-five linear
matrices would lose this action. The affine models used in MF.2
and FP.41 retain it.

\paragraph{FC.3 (Proved): the characteristic elliptic cover and the scope of its geometric interpretation.}
The geometric Ferrari identification of Nakamura, Tsunogai, and
Yasuda (2010, Sections 2.9 and 2.15) identifies $K$ with the
geometric fundamental group of a once-punctured elliptic curve
$E\setminus\{O\}$. Under this identification, $K'K^2$ is the
kernel of the mod-two abelianization map. Its associated connected
degree-four cover is exactly
\[
       [2]:E\setminus E[2]\longrightarrow E\setminus\{O\}.
 \tag{FC.3a}
\]
Indeed the proper multiplication-by-two covering corresponds to
$2H_1(E,\mathbb Z)$, and filling the single puncture identifies
$H_1(E\setminus\{O\},\mathbb Z/2)$ with
$H_1(E,\mathbb Z/2)$. This proves the asserted equality of covering
subgroups. A path to a chosen lift of the basepoint changes only
the inner representative of the identification.

The resulting geometric exact sequence is
\[
 0\longrightarrow
 \mathbb F_3[E[2]]\Big/\mathbb F_3\sum_{P\in E[2]}P
 \longrightarrow H_1(E\setminus E[2],\mathbb F_3)
 \longrightarrow H_1(E,\mathbb F_3)\longrightarrow0.
 \tag{FC.3b}
\]
It has dimensions $3,5,2$. The deck group $Q_0=E[2]$
acts trivially on the proper quotient and regularly on the four
punctures, so averaging over $Q_0$ gives precisely the canonical
splitting (FC.2c). In this interpretation the involution with
class $(1,1,1)\in Q$ fixes all four punctures. The individual
$t_i$ represent that involution composed with the three nonzero
two-torsion translations; they do not each fix all punctures.
This accounts for the boundary weights being two rather than one.

This geometric identification is a statement about a characteristic
cover and its geometric homology. It supplies no unproved
all-$\widehat{GT}$ comparison of the special-fibre paths in
Nakamura--Tsunogai--Yasuda, Proposition 4.13. In particular it does
not identify the cubic frame character $\kappa_d$ with the
quadratic character $\nu_4$.

\paragraph{FC.4 (Audit conclusion): the first stage is insufficient, and the next stage is genuinely new.}
The affine model of FP.41 and MF.2 is consistent with this
independent computation. Its five pure pentagon inputs act on
$B$ by translations, an abelian group; every commutator word
$f\in\widehat F'$ therefore has trivial boundary evaluation.
On $P$, the saved common affine conjugation identifies the
target with the previously audited order-$216$ quotient, with
the same marked five inputs. Thus the stated first-stage
collapse follows; it is not an inference from the dimensions
or an abstract isomorphism of unmarked finite groups.

The mixed class-two construction in MF.3 retains the six-dimensional
bracket module $B\otimes P$. Its definition as the quotient of
the full class-two exponent-three group by
$\bigwedge^2B\oplus\bigwedge^2P$ is structurally legitimate:
the canonical $L$-invariant splitting above makes this discarded
subspace invariant. The resulting target retains information
absent from the first homology stage. Whether its pentagon
equation imposes a new relation, and whether that relation
identifies $\nu_4$ with $\kappa_d\bmod2$, remain separate
questions; neither conclusion follows from this construction.

\paragraph{Primary source.}
Nakamura, H., Tsunogai, H., \& Yasuda, S. (2010).
Harmonic and equianharmonic equations in the
Grothendieck--Teichm\"uller group. III.
\emph{Journal of the Institute of Mathematics of Jussieu, 9}(2),
431--448. \texttt{https://doi.org/10.1017/S1474748009000115}.
The exact Ferrari commutator relation is in Section 2.9 and
the universal punctured elliptic-curve interpretation in
Section 2.15. The group and cube calculations above have
independent proofs and a separate executable certificate.
